\documentclass[11pt]{article}
\usepackage[T1]{fontenc}
\usepackage[margin=1in]{geometry}
\usepackage{amsmath,amssymb}
\usepackage{microtype}
\usepackage{xcolor}
\usepackage[colorlinks=true,linkcolor=blue!55!black,urlcolor=blue!65!black]{hyperref}

\newcommand{\R}{\mathbb R}
\newcommand{\norm}[1]{\lVert#1\rVert}
\newcommand{\sig}{\sigma}

\title{The Semyanistyi weighted-norm question in arXiv:1207.5180\\
was answered by arXiv:1210.5414}
\author{Literature-resolution packet for arXiv:1207.5180}
\date{13 August 2026}

\begin{document}
\maketitle

\begin{center}
\fbox{\begin{minipage}{0.91\textwidth}
\textbf{Status: literature already answered.}
Rubin's separate follow-up arXiv:1210.5414 proves the sharp weighted norm of
the Semyanistyi fractional $k$-plane integral.  Its formula has exactly the
conjectured limit as $\alpha\to0$.  Banach-space adjoint duality gives the
corresponding sharp result and norm limit for the dual operators.
\end{minipage}}
\end{center}

\section{The source question}

Section 5, item 3 (printed page 14) of arXiv:1207.5180 considers
\[
 (R_k^\alpha f)(\tau)=\frac{1}{\gamma_{n-k}(\alpha)}
 \int_{\R^n}f(x)|x-\tau|^{\alpha+k-n}\,dx,
 \qquad
 \gamma_{n-k}(\alpha)=
 \frac{2^\alpha\pi^{(n-k)/2}\Gamma(\alpha/2)}
 {\Gamma((n-k-\alpha)/2)}.
\]
It conjectures that the paper's sharp weighted-$L^p$ method extends to these
operators and their duals, and expects
\[
 \lim_{\alpha\to0}\norm{R_k^\alpha}=\norm{R_k}
\]
on the relevant weighted spaces.  This is the complete mathematical content
of item 3; items 1, 2, and 4 are different questions and are not classified
by this packet.

\section{The direct later theorem}

Rubin's arXiv:1210.5414, submitted three months later, is titled
\emph{Weighted norm estimates for the Semyanistyi fractional integrals and
Radon transforms}; it cites arXiv:1207.5180 explicitly.  Its Theorem 3.2
(printed pages 8--9) states the sharp result for every $1\le p\le\infty$ and
$\alpha>0$.  With
\[
 \nu=\mu-\alpha-k/p',
 \qquad \alpha+k-n/p<\mu<n/p',
\]
one has
\[
 \norm{R_k^\alpha f}_{p,\nu}^{\sim}
 \le c_{\alpha,k}\norm f_{p,\mu},
 \qquad c_{\alpha,k}=\norm{R_k^\alpha},
\]
where
\begin{align*}
c_{\alpha,k}
={}&2^{-\alpha}\pi^{k/2}
 \left(\frac{\sig_{n-k-1}}{\sig_{n-1}}\right)^{1/p}\\
&\times
\frac{\Gamma((n/p'-\mu)/2)
      \Gamma((\mu+n/p-k-\alpha)/2)}
     {\Gamma((\mu+n/p)/2)
      \Gamma((n/p'-\mu+\alpha)/2)}.
\end{align*}
Thus the later theorem supplies boundedness, the admissible power weights,
and the exact operator norm, rather than only an estimate.  Theorem 3.3
(printed page 12) records the zero-order theorem in the same notation:
\[
 \norm{R_k}=c_k=
 \pi^{k/2}\left(\frac{\sig_{n-k-1}}{\sig_{n-1}}\right)^{1/p}
 \frac{\Gamma((\mu+n/p-k)/2)}
      {\Gamma((\mu+n/p)/2)},
 \quad \nu=\mu-k/p'.
\]

\section{The asserted norm limit}

Fix $p$ and a weight exponent $\mu$ with
\[
 k-n/p<\mu<n/p'.
\]
For all sufficiently small $\alpha>0$, the fractional theorem applies with
the compatible target weight $\nu_\alpha=\mu-\alpha-k/p'$.  Continuity of
$2^{-\alpha}$ and of $\Gamma$ away from its poles gives
\begin{align*}
\lim_{\alpha\downarrow0}c_{\alpha,k}
={}&\pi^{k/2}\left(\frac{\sig_{n-k-1}}{\sig_{n-1}}\right)^{1/p}
\frac{\Gamma((n/p'-\mu)/2)
      \Gamma((\mu+n/p-k)/2)}
     {\Gamma((\mu+n/p)/2)
      \Gamma((n/p'-\mu)/2)}\\
={}&c_k=\norm{R_k}.
\end{align*}
This is precisely the expected limit, with the naturally varying target
weight tending to $\mu-k/p'$.  The upper restriction $\mu<n/p'$ is needed
for the positive-order family; hence the limit is asserted on the common
interior of the fractional weighted spaces, not at a parameter where
$R_k^\alpha$ is unbounded.

\section{Dual operators}

The introduction of arXiv:1210.5414 explicitly says that the same approach
applies to the dual Radon transforms and dual Semyanistyi integrals.  More
decisively, the dual conclusion follows immediately from its sharp primal
theorem.  Under the unweighted pairing of the two underlying measure spaces,
\[
 (R_k^\alpha)^*:
 L^q_\beta(V_{n,n-k}\times\R^{n-k})
 \longrightarrow L^q_\delta(\R^n),
 \qquad \delta=\beta-\alpha-k/q,
\]
is obtained by applying Theorem 3.2 with $p=q'$, $\mu=-\delta$, and
$\nu=-\beta$.  The admissible range becomes
\[
 \alpha-(n-k)/q<\beta<(n-k)/q',
\]
and $\norm{(R_k^\alpha)^*}=\norm{R_k^\alpha}$ for the corresponding paired
weights.  Substitution into the displayed gamma formula gives the sharp dual
constant.  Letting $\alpha\downarrow0$ cancels the same gamma pair and yields
the sharp dual $R_k^*$ norm in Theorem 1.1 of arXiv:1207.5180.  Therefore the
word ``duals'' in the source question creates no unresolved residue.

\section{Conclusion and audit}

The source's item 3 and the later paper match in operator, weighted spaces,
sharpness, and limiting statement.  Because the supporting theorem was
published as a separate paper by the same author and cites the source paper,
this is a direct literature answer, not an independently inferred partial
analogy.

\paragraph{Source audit.}
The official arXiv PDFs and source files of 1207.5180 and 1210.5414 were
inspected.  The source question is Section 5, item 3, printed page 14.  The
decisive later results are Theorems 3.2 and 3.3, printed pages 8--9 and 12.
The target was absent from the run's cheap solution indexes before this
packet.

\paragraph{Human review.}
Check the endpoint convention $\alpha\downarrow0$ and the reparameterization
in the one-line adjoint argument.  The primal theorem identification and
gamma-factor cancellation are literal.

\begin{thebibliography}{9}
\bibitem{R12a}
B.~Rubin,
\emph{Weighted norm inequalities for $k$-plane transforms},
arXiv:1207.5180.

\bibitem{R12b}
B.~Rubin,
\emph{Weighted norm estimates for the Semyanistyi fractional integrals and
Radon transforms}, arXiv:1210.5414.
\end{thebibliography}

\end{document}
